\section*{Conflict of interest}

All authors are affiliated with \policyengine{}, the nonprofit organization
that develops \populace{}, \populacefit{}, and \microimpute{}, the software
evaluated in this paper. The authors have no other competing interests.

\section*{Funding}

\todo{Funding statement.}

\section*{Data and code availability}

The paper and experiment code are at
\url{https://github.com/PolicyEngine/imputation-paper}; every results table
regenerates from committed run configurations via the repository's command
line. The estimator is implemented in \populacefit{}, part of \populace{},
open source at \url{https://github.com/PolicyEngine/populace}. Baseline
implementations are \microimpute{}
(\url{https://github.com/PolicyEngine/microimpute}) and \pystatmatch{}
(\url{https://github.com/CosilicoAI/py-statmatch}). The reported sweeps ran
against \populacefit{} at populace commit \texttt{71b5a83},
\microimpute{} 1.1.2, and \pystatmatch{} at commit \texttt{094a242}; the
\scf{} 2022 summary extract (SHA-256 beginning \texttt{3bb4d890}) and the raw
\texttt{cps\_2023.h5} (SHA-256 beginning \texttt{b7b68ee0}) are the pinned
survey inputs, and each run directory's manifest records rows, caps, seeds,
and methods.
